% !TeX root = ../main.tex
\chapter{Theoretische Grundlagen}
\label{ch:theory}
\section{Domänenspezifische Grundlagen: SAP Variantenkonfiguration}
\subsection{Einführung in die logische Architektur der Variantenkonfiguration}

In der modernen Fertigungsindustrie erfordert die Abbildung kundenindividueller Produkte hochflexible 
Systemstrukturen. Die grundlegende informationstechnische Idee der Variantenkonfiguration besteht darin, ein 
Produkt durch eine Vielzahl möglicher Ausprägungen zu definieren, welche durch ein iteratives Regelwerk logisch 
miteinander verknüpft sind, um komplexe Geschäftsprozesse strukturiert abzubilden. Da die Ausmultiplikation aller 
möglichen Kombinationsmöglichkeiten die Anzahl der Varianten schnell in die Millionen oder in extremen Fällen der 
Anlagenfertigung bis in den Bereich von Quadrillionen ($10^{24}$) treiben kann, ist die Anlage diskreter Stammdaten 
für jede einzelne Ausprägung in der Praxis unmöglich (zitat). 

Als systemtechnische Antwort auf diese kombinatorische Explosion fungiert im SAP-System der sogenannte 
\textit{konfigurierbare Materialstamm} (KMAT). 
Dieser dient als universaler Platzhalter, unter dem die gesamte Variantenvielfalt eines Produkts 
informationstechnisch zusammengefasst wird (zitat). 
Um diese Vielfalt in einer sogenannten Bewertungsoberfläche abzubilden, greift die SAP-Architektur auf 
drei zentrale Strukturkomponenten zurück (zitat):

\begin{itemize}
    \item \textbf{Merkmale (Characteristics):} Sie definieren die spezifischen Eigenschaften eines Produkts 
    (wie beispielsweise die Lackierung oder die Traglast eines Gabelstaplers).
    \item \textbf{Variantenklassen:} Diese dienen der logischen Bündelung der Merkmale und werden direkt dem 
    konfigurierbaren Materialstamm zugeordnet.
    \item \textbf{Konfigurationsprofil:} Dieses fungiert als übergeordnetes Steuerungselement, das die Stammdaten 
    zusammenhält und den Konfigurationsprozess maßgeblich regelt.
\end{itemize}

Ein zentrales Paradigma der SAP-Variantenkonfiguration zur Bewältigung dieser Komplexität ist das sogenannte 
Maximal-Prinzip. Anstatt ein Produkt aus einzelnen Modulen additiv zu einem 100\,\%-Modell zusammenzusetzen, 
basiert die Architektur auf einem subtraktiven 150\,\%-Modell. Hierfür wird eine 
\textit{konfigurierbare Maximalstückliste} (Super BOM) eingesetzt. Diese umfasst sämtliche Bauteile und 
Komponenten, die für jedwede theoretisch mögliche Variante des Produkts jemals benötigt werden könnten 
(zitat). Während des Konfigurationsprozesses können dieser Stückliste keine neuen 
Elemente hinzugefügt, sondern lediglich nicht benötigte Komponenten herausgefiltert werden.

Die logische Brücke zwischen der kaufmännischen Konfiguration und der technischen Fertigung bildet das 
\textit{Beziehungswissen} (Object Dependencies). Metaphorisch gesprochen übersetzt dieses Regelwerk das 
kundenseitige „Wünschen“ in der Bewertungsoberfläche in das fertigungsseitige „Erhalten“ in Form der ausgedünnten 
Stückliste (zitat). Konkret wird dies durch Steuerungsmechanismen wie 
\textit{Auswahlbedingungen} (Selection Conditions) realisiert. Wenn ein Kunde beispielsweise aus einem Angebot von 
sechs Motoren einen spezifischen Antrieb wählt, determinieren die Auswahlbedingungen des Beziehungswissens, 
dass die fünf nicht gewählten Motoren aus der finalen Fertigungsstückliste eliminiert werden 
(zitat). 

Die Zusammenhänge dieser fundamentalen Architekturkomponenten – vom konfigurierbaren Materialstamm über die 
Bewertungsoberfläche bis hin zur Maximalstückliste und dem steuernden Beziehungswissen – sind in Abbildung 
\ref{fig:variantenkonfig} schematisch zusammengefasst.

(Hier kommt die Grafik rein)

\subsection{Die klassische Variantenkonfiguration (LO-VC) im Kontext der Systemtransformation}
\label{sec:lovc_context}

Um die technologische Entwicklung der Variantenkonfiguration zu verstehen, ist zunächst eine Einordnung 
in die Evolution der SAP-Systemlandschaften erforderlich. Während das bisherige Kernprodukt 
\textit{SAP ERP} (oft auch als SAP R/3 bezeichnet) über Jahrzehnte den Standard darstellte, markiert 
\textit{SAP S/4HANA} eine fundamentale Neuausrichtung der Softwarearchitektur, insbesondere durch die 
Nutzung der In-Memory-Datenbanktechnologie \parencite[vgl.][S. 28]{blumohrAdvancedVariantConfiguration2023}. 
In diesem Zuge wurde für die Variantenkonfiguration ein neues technologisches Fundament geschaffen.

Der Begriff \textit{Logistics Variant Configuration} (LO-VC) dient heute primär als Retronym, um die 
klassische Art der Konfiguration von der neuen \textit{Advanced Variant Configuration} (AVC) abzugrenzen 
\parencite[vgl.][S. 32]{blumohrAdvancedVariantConfiguration2023}. Obwohl LO-VC in aktuellen S/4HANA-Installationen 
(On-Premise und Private Cloud) aus Gründen der Abwärtskompatibilität weiterhin verfügbar ist, gilt sie 
technologisch als veralteter Standard. 

Wie bereits in den Grundlagen erläutert, basiert LO-VC auf einem prozeduralen Paradigma, das historisch 
bedingt oft durch einen „Learning-by-Doing“-Ansatz in den Unternehmen gewachsen ist. Dies führte dazu, 
dass Modelle häufig unnötig komplex durch den Einsatz von Prozeduren und Aktionen anstatt 
leistungsfähigerer Constraints gestaltet wurden \parencite[vgl.][S. 58]{blumohrAdvancedVariantConfiguration2023}. Der Wechsel zu 
AVC wird in der Literatur daher nicht nur als technisches Upgrade, sondern als strategische Notwendigkeit 
beschrieben, die durch folgende Kernaspekte begründet wird:

\begin{itemize}
    \item \textbf{Rückkehr zum Standard und Vereinfachung:} AVC stellt den künftigen Standard dar. Die 
    Migration bietet Unternehmen die Chance, historisch gewachsene, schwer wartbare „Spaghetti-Logiken“ 
    zu dekonstruieren und durch eine moderne, leichter wiederverwendbare Modellierung zu ersetzen 
    \parencite[vgl.][S. 57\,f.]{blumohrAdvancedVariantConfiguration2023}.
    \item \textbf{Abbildung moderner Geschäftsmodelle:} Klassische LO-VC-Strukturen sind primär auf den 
    reinen Verkauf von Produkten (SD-Kontext) ausgelegt. Moderne Strategien wie das 
    \textit{Solution Bundling} (Kombination aus Produkt, Dienstleistung und Abonnements) lassen sich im 
    Standard nur über AVC und die damit verknüpften neuen Objektträger wie die \textit{Solution Quote} 
    abbilden \parencite[vgl.][S. 59]{blumohrAdvancedVariantConfiguration2023}.
    \item \textbf{Technologische Überlegenheit:} AVC nutzt mit \textit{Gecode} eine moderne 
    Constraint-Solving-Engine, die im Gegensatz zur sequenziellen LO-VC-Logik analytische Prinzipien 
    verfolgt und damit eine deutlich höhere Performance und Flexibilität bietet 
    \parencite[vgl.][S. 59]{blumohrAdvancedVariantConfiguration2023}.
    \item \textbf{Zukunftsfähigkeit (KI und Analytics):} Innovative Technologien wie 
    \textit{Artificial Intelligence} (AI) oder \textit{Machine Learning} (ML) sowie moderne 
    Auswertungsmethoden via \textit{CDS-Views} werden nativ für AVC entwickelt. Eine nachträgliche 
    Integration in die veraltete LO-VC-Architektur gleicht lediglich einer „Renovierung“, die nie den 
    Leistungsumfang einer Neuentwicklung erreichen kann \parencite[vgl.][S. 61]{blumohrAdvancedVariantConfiguration2023}.
    \item \textbf{Human-Resource-Aspekt:} Angesichts des Fachkräftemangels ist die Arbeit mit moderner 
    Softwarearchitektur ein entscheidender Faktor bei der Gewinnung von Talenten. Die Wartung einer über 
    25 Jahre alten Softwarelösung ist im Vergleich zur Vorreiterrolle in einer intelligenten 
    Fertigungsumgebung auf Basis aktueller Technologie weniger attraktiv für qualifizierten Nachwuchs 
    \parencite[vgl.][S. 60]{blumohrAdvancedVariantConfiguration2023}.
\end{itemize}

Zusammenfassend lässt sich festhalten, dass der Verbleib in der LO-VC-Welt zwar kurzfristig 
Transformationsaufwände vermeidet, langfristig jedoch das Risiko erhöht, den Anschluss an technologische 
Innovationen zu verlieren und den Aufwand für eine spätere Migration durch das wachsende technologische 
Delta exponentiell zu vergrößern \parencite[vgl.][S. 61]{blumohrAdvancedVariantConfiguration2023}.

\subsection{Die moderne Architektur (AVC) und das Constraint-Paradigma}
\label{sec:avc_constraint}

Die Antwort auf die konzeptionellen Schwächen und technischen Limitierungen der klassischen 
Variantenkonfiguration bildet die Einführung der \textit{Advanced Variant Configuration} (AVC) in SAP 
S/4HANA. Um den stetig wachsenden Anforderungen an Performance und Modellierungskomplexität 
(wie beispielsweise in den \textit{Make-to-Order}- oder \textit{Engineer-to-Order}-Szenarien) gerecht zu 
werden, vollzieht SAP mit AVC nicht nur ein syntaktisches Update, sondern einen fundamentalen 
informationstechnischen Paradigmenwechsel unter der Haube der Konfigurations-Engine 
\parencite[vgl.][S. 35]{blumohrAdvancedVariantConfiguration2023}.

Während das bisherige LO-VC-System, wie in Kapitel \ref{sec:lovc_context} beschrieben, auf einer 
iterativen, imperativen Abarbeitung von Befehlen und Prozeduren basierte, nutzt AVC einen grundlegend 
deklarativen Lösungsansatz. Das technologische Herzstück der AVC-Architektur bildet ein neu integrierter 
\textit{Constraint Solver} namens \textit{Gecode}, eine ursprünglich quelloffene Softwarebibliothek, 
die in Zusammenarbeit mit dem Fraunhofer-Institut ITWM tief in den SAP-Standard eingebettet und für 
komplexe Produktkonfigurationen erweitert wurde \parencite[vgl.][S. 35]{blumohrAdvancedVariantConfiguration2023}.

Der Einsatz dieser Engine überführt den Konfigurationsprozess mathematisch in ein sogenanntes 
\textit{Constraint Satisfaction Problem} (CSP). Bei einem CSP geht es aus informatischer Sicht nicht mehr 
darum, eine Liste von Instruktionen nacheinander abzuarbeiten, um einen Endzustand zu berechnen. 
Stattdessen wird das gesamte Modell als ein Netzwerk aus Variablen (den SAP-Merkmalen) und den 
zugehörigen Wertebereichen (den Domänen) definiert. Das Beziehungswissen wird in diesem Paradigma als 
harte Bedingungen (\textit{Constraints}) formuliert, die festlegen, welche Merkmalskombinationen 
zulässig sind und welche sich gegenseitig ausschließen. Ziel des CSP-Solvers ist es, eine gültige 
Belegung für alle Variablen zu finden, bei der kein einziger Constraint verletzt wird. 

Dieser deklarative Ansatz bringt tiefgreifende Auswirkungen auf die Benutzerführung und die 
Systemstabilität mit sich. Sobald ein Anwender während der Konfiguration in der Benutzeroberfläche einen 
Wert für ein bestimmtes Merkmal wählt, schränkt die Gecode-Engine die Domänen (Wertebereiche) aller 
anderen abhängigen Merkmale in Echtzeit ein \parencite[vgl.][S. 35]{blumohrAdvancedVariantConfiguration2023}. 
Anstatt einen Fehler erst nach der sequenziellen Abarbeitung einer Prozedur auszuwerfen, eliminiert der 
Solver unmögliche Auswahlmöglichkeiten präventiv aus dem Suchraum. Dies verhindert effektiv die 
Entstehung inkonsistenter Konfigurationen.

Darüber hinaus bietet AVC erweiterte Simulations- und Analysemöglichkeiten, wie beispielsweise den 
\textit{Dependency Trace}. Dieser ermöglicht es Konstrukteuren und Modellierern, in einem 
\textit{Inspector Panel} genau nachzuvollziehen, welches Constraint für die Einschränkung eines 
bestimmten Merkmals verantwortlich ist \parencite[vgl.][S. 37]{blumohrAdvancedVariantConfiguration2023}. 

Zusammenfassend markiert AVC durch die Gecode-Engine den Wechsel von der Frage „\textit{Wie} berechne 
ich die Variante Schritt für Schritt?“ hin zu der Frage „\textit{Welche} globalen Bedingungen müssen für 
ein gültiges Endprodukt erfüllt sein?“. Genau dieser Paradigmenwechsel vom Imperativen zum Deklarativen 
stellt Unternehmen bei der Systemmigration vor die immense Herausforderung, historisch gewachsene, 
prozedurale Altlasten semantisch zu dekonstruieren und in das neue, netzwerkbasierte Constraint-Paradigma 
zu transformieren.

\section{Generative künstliche Intelligenz im Software Engineering}
\subsection{Large Language Models für Quellcode}
\label{subsec:llms_fuer_code}

In den vergangenen Jahren haben Systeme der Künstlichen Intelligenz (KI), insbesondere sogenannte Large 
Language Models (LLMs), einen rasanten und tiefgreifenden technologischen Aufstieg erlebt. Was als 
Werkzeug für rein statistische Textvervollständigungen begann, hat sich zu mächtigen kognitiven 
Architekturen entwickelt. Diese Modelle dringen zunehmend in den Bereich des Software Engineerings vor 
und verändern die Art und Weise, wie Software entworfen, analysiert und transformiert wird, grundlegend. 
Um das Potenzial dieser Technologie für automatisierte Migrationsprozesse zu verstehen, ist ein 
grundlegendes Verständnis der zugrundeliegenden mathematisch-technischen Prinzipien sowie der Methoden 
zu deren Steuerung erforderlich.

\subsubsection{Die Transformer-Architektur und der Aufmerksamkeitsmechanismus}
Die heutige Generation marktbeherrschender Sprachmodelle basiert fast ausnahmslos auf der wegweisenden 
Transformer-Architektur, die durch \cite{vaswaniAttentionAllYou2017} eingeführt wurde. Der entscheidende 
informationstechnische Durchbruch des Transformers liegt in der vollständigen Abkehr von früheren 
rekurrenten neuronalen Netzen (RNNs) oder Long Short-Term Memory-Netzwerken (LSTMs). Während diese 
älteren Architekturen Eingabesequenzen – wie natürlichen Text oder Quellcode – streng sequenziell, 
mithin Zeichen für Zeichen oder Zeile für Zeile verarbeiten mussten, bricht der Transformer dieses 
lineare Paradigma auf.

Dies gelingt durch die Einführung des sogenannten \textit{Self-Attention-Mechanismus} 
(Selbstaufmerksamkeit). Dieser Mechanismus erlaubt es dem Modell, alle Elemente (Token) einer Sequenz 
zeitgleich und vollständig parallel zu verarbeiten. Mathematisch berechnet das Netzwerk für jedes Token 
innerhalb eines Kontextes einen Gewichtungsfaktor relativ zu allen anderen Token der Sequenz. Dadurch 
wird bestimmt, wie viel „Aufmerksamkeit“ einem Wort oder Code-Snippet im Verhältnis zum Rest des Textes 
beigemessen werden muss. 

Im Kontext des Software Engineerings löst diese Architektur das fundamentale Problem weitreichender 
Abhängigkeiten (\textit{Long-Range Dependencies}). Wenn beispielsweise eine Variable am Anfang einer 
umfangreichen Code-Datei deklariert und hunderte Zeilen später in einer komplexen logischen Bedingung 
aufgerufen wird, verkürzt der Transformer den informationellen Pfad im Netzwerk auf eine konstante Länge. 
Der Self-Attention-Mechanismus berechnet die direkte strukturelle Verknüpfung zwischen Deklaration und 
Aufruf unabhängig von ihrer physischen Distanz im Dokument. Quellcode wird somit nicht als linearer Text, 
sondern als ein Netz simultan interagierender logischer Einheiten interpretiert.

\subsubsection{In-Context Learning und Few-Shot Prompting}
Die Fähigkeit, tiefe semantische Strukturen zu erfassen, wirft die Frage auf, wie diese Modelle ohne 
immensen Rechenaufwand für spezifische, proprietäre Aufgabensteuerungen eingesetzt werden können. 
Historisch erforderte die Anpassung eines neuronalen Netzes an eine neue Domäne ein ressourcenintensives 
\textit{Fine-Tuning}, bei dem die internen Gewichte des Modells durch überwachtes Lernen angepasst 
wurden. Moderne LLMs weisen jedoch ab einer kritischen Skalierungsgröße sogenannte emergente Fähigkeiten 
auf, die dieses Paradigma verschieben.

Wie durch \cite{brownLanguageModelsAre2020} nachgewiesen wurde, verhalten sich ausreichend große Sprachmodelle 
weitgehend aufgabenunabhängig (\textit{task-agnostisch}). Anstatt die Modellparameter dauerhaft zu 
modifizieren, findet das Lernen direkt zur Laufzeit innerhalb des Eingabefensters statt – ein Konzept, 
das als \textit{In-Context Learning} (ICL) bezeichnet wird. Das Modell nutzt sein im Pre-Training 
erworbenes, generalistisches Strukturwissen, um eine neue Aufgabe fließend und unmittelbar aus den im 
Prompt gegebenen Mustern abzuleiten.

Eine der mächtigsten Methoden des In-Context Learnings ist das \textit{Few-Shot Prompting}. Hierbei 
werden dem Modell im Eingabetext neben der eigentlichen Instruktion einige wenige konkrete 
Demonstrationen (Musterbeispiele) übergeben, die eine Eingabe und die dazugehörige korrekte Zielausgabe 
abbilden. Ohne dass ein einziges Gradientenupdate im neuronalen Netz stattfindet, adaptiert das LLM die 
logische Analogie zwischen den Beispielen und wendet diese fehlerfrei auf die neue, ungesehene 
Zielaufgabe an. Für komplexe Domänen bietet dies den Vorteil, dass Modelle flexibel und agil gesteuert 
werden können, solange der Prompt die notwendigen logischen Leitplanken bereitstellt.

\subsubsection{Das Paradoxon der Code-Übersetzung}
Für den konkreten Einsatz generativer KI im Software Engineering liefert die umfassende Systematisierung 
von \cite{zhengSurveyLargeLanguage2023} eine wesentliche differenzierende Erkenntnis bezüglich der Modellauswahl. 
Entgegen der intuitiven Annahme, dass für Programmieraufgaben spezialisierte Modelle stets vorzuziehen 
sind, zeigt sich in empirischen Studien ein zweigeteiltes Bild.

Auf der einen Seite existieren hochspezialisierte, verhältnismäßig kleine Sprachmodelle 
(\textit{Code LLMs}), die exklusiv auf riesigen Quellcode-Korpora trainiert wurden. Beim Agieren auf 
einer „grünen Wiese“ – mithin bei der isolierten Generierung von Standardfunktionen oder der 
automatisierten Erstellung von Unit-Tests – übertreffen diese dedizierten Code-Modelle die gigantischen, 
allgemeinsprachlichen Modelle oft deutlich. Sie weisen eine hohe syntaktische Präzision auf und sind in 
ihrer Ausführung hochgradig effizient.

Auf der anderen Seite verkehrt sich dieses Leistungsverhältnis ins Gegenteil, sobald die Aufgabe in den 
Bereich der \textit{Code-Übersetzung} (\textit{Code Translation}) fällt. Bei der Migration bestehender 
Logik von einer Quellsprache oder einem Paradigma in eine andere Zielarchitektur dominieren führende, 
generalistische Großmodelle (wie beispielsweise GPT-4) die spezialisierten Code-Modelle signifikant. 

Die wissenschaftliche Begründung hierfür liegt in der Natur der Problemstellung: Eine Software-Migration 
ist kein rein syntaktisches Substitutionsproblem, bei dem lediglich Schlüsselwörter ausgetauscht werden. 
Es handelt sich vielmehr um ein komplexes Abstraktions- und Logikproblem. Das Modell muss den 
semantischen Kern, die implizite Absicht und die vernetzten Bedingungen des Altsystems vollständig 
durchdringen, diese in eine sprachunabhängige logische Repräsentation abstrahieren und anschließend 
unter den restriktiven Bedingungen des Zielparadigmas neu formulieren. Für derart vielschichtige 
logische Schlussfolgerungen (\textit{Reasoning}) ist das breite Weltwissen und die schiere kognitive 
Kapazität der großen, generalistischen Modelle unabdingbar, während rein syntaxfokussierte Code-Modelle 
an den abstrakten Paradigmenwechseln scheitern.

\subsection{Limitationen monolithischer Ansätze}
\label{subsec:limitationen_monolithischer_ansaetze}

Obgleich die im vorherigen Abschnitt beschriebenen Fähigkeiten moderner Großmodelle zur semantischen 
Erfassung und Übersetzung von Quellcode beachtlich sind, stößt der praktische Einsatz in industriellen 
Software-Engineering-Projekten an fundamentale Grenzen. In der Praxis und der aktuellen Forschung werden 
Ansätze, bei denen eine komplexe Aufgabe in einer einzigen Interaktionsschleife – mithin über einen 
einzelnen, isolierten Eingabeaufforderungs-Befehl (\textit{Single-Prompt} oder 
\textit{One-Shot}-Verfahren) – gelöst werden soll, als \textit{monolithische Ansätze} bezeichnet. 
Um die Notwendigkeit fortgeschrittener, verteilter Architekturen zu begründen, müssen die 
systemimmanenten Schwachstellen dieser monolithischen Herangehensweise auf zwei Ebenen analysiert 
werden: der algorithmischen Komplexität der Logikketten und der strukturellen Komplexität systemweiter 
Abhängigkeiten.

\subsubsection{Die Komplexitätsfalle und das probabilistische Paradigma}
Die primäre algorithmische Limitierung monolithischer Prompts bei logisch anspruchsvollen Aufgaben 
wurde bereits frühzeitig durch OpenAI im Rahmen der Vorstellung des Codex-Modells empirisch dokumentiert. 
Wie \cite{chenEvaluatingLargeLanguage2021} nachweisen, sind LLMs hervorragend in der Lage, isolierte, in sich 
geschlossene Programmieraufgaben zu lösen, sofern die Spezifikation klar umrissen ist. 

Das Leistungsvermögen bricht jedoch mathematisch nachweisbar ein, sobald eine Spezifikation lange 
Ketten sequenzieller oder verschachtelter Operationen verlangt (\textit{Chaining-Problem}). Die Autoren 
zeigen, dass die Erfolgsquote des Modells mit zunehmender Anzahl logisch verknüpfter Bausteine 
innerhalb der Anforderung exponentiell abfällt \cite{chenEvaluatingLargeLanguage2021}. Für die Migration historisch 
gewachsener SAP-Regelwerke (wie dem klassischen \textit{Variant Configuration}-Modell LOVC) ist diese 
Erkenntnis von kritischer Bedeutung: Solche Altsysteme bestehen typischerweise aus hochgradig 
verschachtelten, prozeduralen Abhängigkeiten und logischen Verzweigungen. Ein monolithischer Prompt 
zwingt das Modell, diese gesamte kognitive Last simultan zu verarbeiten und den Zielcode in einem 
einzigen autoregressiven Durchlauf fehlerfrei zu generieren. Dies überschreitet die inhärenten 
Fähigkeiten zur logischen Vorausschau des Modells, da es Token für Token auf Basis von 
Wahrscheinlichkeiten generiert, ohne eine echte logische Tiefenprüfung vorzunehmen.

Darüber hinaus beleuchtet das Paper eine fundamentale Schwäche des probabilistischen Charakters von 
Sprachmodellen. Die Metrik $pass@1$ – welche die Wahrscheinlichkeit beschreibt, dass ein einzelner 
generierter Code-Entwurf funktional korrekt ist – ist bei komplexen Aufgabenstellungen extrem gering. 
Erst durch das wiederholte Generieren unabhängiger Stichproben ($pass@k$) und eine anschließende 
externe Verifizierung steigt die Erfolgsquote signifikant an \cite{chenEvaluatingLargeLanguage2021}. Ein 
monolithischer Ansatz bietet jedoch keinerlei inhärente Mechanismen für ein solches iteratives 
\textit{Sampling} oder eine automatisierte Qualitätskontrolle, wodurch fehlerhafte Syntax oder logische 
Fehlschlüsse unkorrigiert in das Zielergebnis einfließen.

\subsubsection{Das Scheitern an systemweiten Abhängigkeiten und fehlendem Umgebungs-Feedback}
Während die Arbeit von Chen et al. die Schwächen auf Ebene einzelner Funktionen aufzeigt, skaliert der 
aktuellere \textit{SWE-bench}-Benchmark von \cite{jimenezSwebenchCanLanguage2024} diese Analyse auf das Niveau 
realer, industrieller Software-Repositories. Im Rahmen dieses Frameworks wurden führende Sprachmodelle 
mit echten, komplexen Problemstellungen aus großen Open-Source-Projekten konfrontiert, bei denen 
Änderungen nicht isoliert, sondern über Systemgrenzen hinweg durchgeführt werden mussten.

Die Ergebnisse dieses Benchmarks offenbaren ein klares Defizit monolithischer Großmodelle: Selbst 
hochgradig leistungsfähige, kommerzielle Flaggschiff-Modelle wie GPT-4 oder Claude waren im 
monolithischen \textit{Single-Prompt}-Setting überfordert und konnten weit unter $5\%$ der realen 
Software-Probleme erfolgreich lösen \cite{jimenezSwebenchCanLanguage2024}. Die Autoren identifizieren zwei 
Hauptursachen für dieses drastische Scheitern:

\begin{enumerate}
    \item \textbf{Mangelndes Verständnis globaler Abhängigkeiten:} Monolithische LLMs scheitern an der 
    Identifikation von \textit{Cross-File-Dependencies}. Sie sind zwar in der Lage, den lokal 
    übergebenen Code-Kontext zu modifizieren, überblicken jedoch nicht, welche Auswirkungen diese 
    Modifikation auf andere, im Kontext nicht enthaltene Klassen, Schnittstellen oder Datenstrukturen 
    des Gesamtsystems hat. Bei einer SAP AVC-Migration, die auf einem global vernetzten, deklarativen 
    \textit{Constraint}-Netzwerk basiert, führt dieser blinde Fleck unweigerlich zu fatalen 
    Inkonsistenzen im Gesamtsystem.
    \item \textbf{Fehlen einer interaktiven Ausführungsumgebung:} Ein monolithischer Prompt agiert als 
    reine Text-zu-Text-Transformation in einer isolierten Umgebung (Black-Box). Dem Modell fehlt 
    jeglicher Zugriff auf einen Compiler, einen Interpreter oder Testprotokolle. In der realen 
    Softwareentwicklung ist die fehlerfreie Code-Erstellung jedoch ein hochgradig interaktiver, 
    iterativer Prozess. Ohne die Möglichkeit, den generierten Code testweise auszuführen und das daraus 
    resultierende Compiler-Feedback dynamisch zu interpretieren, ist das Modell gezwungen, die 
    syntaktische Korrektheit blind zu erraten.
\end{enumerate}

Zusammenfassend lässt sich konstatieren, dass monolithische KI-Ansätze aufgrund der exponentiellen 
Fehleranfälligkeit bei langen Logikketten, des Mangels an systemweiter Kontextverarbeitung und der 
Isolation von echten Ausführungsumgebungen für komplexe Architektur-Transformationen ungeeignet sind. 
Diese wissenschaftlich dokumentierte Forschungslücke erfordert den Übergang von isolierten Prompts hin 
zu dynamischen, interaktiven Multi-Agenten-Systemen, die in der Lage sind, Aufgaben modular aufzuteilen 
und durch kontinuierliche Feedbackschleifen zu validieren.